% questions.tex —— 1958 年·湖南卷 · 题目内容片段（无答案）
% 由 tools/generate_projects.py 依来源总账自 1958/1958_湖南_数学真题_试题.md 生成，勿手改（改源 MD 后重跑）


\begin{question}
\gkti{1}不用除法证明 $(2+x)^{2n-1}+x^{2n-1}+x+1$ 能被 $x+1$ 整除（$n$ 为正整数）。

\end{question}
\begin{question}
\gkti{2}不解方程 $6x^2-x-2=0$，求作一个二次方程，使它的根是原方程的根的倒数。

\end{question}
\begin{question}
\gkti{3}求证 $\sin 75^\circ=\dfrac14\left(\sqrt6+\sqrt2\right)$。

\end{question}
\begin{question}
\gkti{4}证明：设三角形的三个角成等差数列，最大角与最小角相差 $30^\circ$，则其三边之比顺次等于 $2\sqrt2:2\sqrt3:(\sqrt6+\sqrt2)$ 之比。

\end{question}
\begin{question}
\gkti{5}一矩形内接于圆，它的面积是 $48\,\text{cm}^2$，圆半径是 $5\,\text{cm}$，求矩形的两邻边。

\end{question}
\begin{question}
\gkti{6}设 $AD,BE$ 是由圆的直径 $AB$ 两端点所引的切线，$DE$ 是过圆上任意一点 $F$ 的切线，此切线与 $AD,BE$ 相交于 $D$ 及 $E$，求证：$AD\cdot BE=OF^2$。

\end{question}
\begin{question}
\gkti{7}正六棱锥底面的边长是 $4\,\text{cm}$，侧面与底面所成二面角等于 $45^\circ$，求这棱锥体的体积。

\end{question}
\begin{question}
\gkti{8}证明二面角的平分面内任意一点到这二面角的两个面的距离相等。

\end{question}